\section{Introduction}
\label{sec:intro}

Motivational Interviewing \citep{miller1991mi} is a counselling style built on a specific
discipline: the therapist elicits the client's own reasons for change rather than supplying
them. Its defining moves are open questions and complex reflections; its defining
prohibitions are unsolicited advice, confrontation, and --- less obviously --- the kind of
warm, unconditional praise that feels supportive but substitutes the therapist's approval
for the client's motivation. This makes MI an unusually good testbed for reward
specification, because the difference between ``sounds like a good session'' and ``is a good
session'' has been written down, operationalised, and used to train human coders
\citep{moyers2016miti}.

It is also a domain where LLM-based evaluation is already established. \citet{yosef2024assessing}
show that an LLM filling in MI session questionnaires from a simulated patient's
perspective produces ratings that are statistically reliable and that separate therapists by
expertise level. That result is what makes the pipeline studied here possible at all: a
small therapist model converses with an LLM-simulated patient, a larger LLM completes
validated MI questionnaires on the transcript, and the resulting score is optimised. No
human is in the loop, so the loop can be run for as many iterations as budget allows.

The question this paper asks is what such a loop actually produces. We train a
Llama-3.2-1B therapist against a \primaryjudge{} patient and grader for ten iterations
under two optimizers that differ in how they explore --- a preference-tree method that
branches candidate replies and trains with DPO \citep{rafailov2023dpo,baruch2025pto}, and a
group-relative method that samples a group of completions per prompt and trains with GRPO
\citep{shao2024deepseekmath} --- holding the reward, the patient population, the base model
and every shared hyperparameter fixed.

On its own terms the loop works, and it works impressively. Both arms improve on all five
global-evaluation instruments with large effect sizes; the preference-tree arm takes MITI
global relational ratings from $3.34$ to $4.61$, crossing the manual's ``good'' threshold,
and \QQ{} --- the training reward --- from $3.00$ to $4.26$ ($\dz=1.43$). Read only through
the grader that supplied the reward, this is a strong positive result about automating
therapist training.

We report three findings that complicate it.

\paragraph{The behavioural signature is affirmation replacing inquiry.} Both optimizers
drift the same way. Therapist turns grow $2.3$--$3.4\times$ longer, affirmations per turn
rise more than fivefold, MI-inconsistent behaviour rises $2.3\times$ (PTO) to $4.0\times$
(GRPO), and questions collapse. The collapse is visible most sharply in a disagreement
between two ways of counting it: a deterministic regex counting question marks per therapist
turn falls $0.83\!\rightarrow\!0.15$ in the GRPO arm, while the oracle's own MITI question
code falls only $0.45\!\rightarrow\!0.32$ over the same conversations. The two measures
start out consistent and end up inverted --- the grader continues to credit inquiry after
the literal questions are gone (\S\ref{sec:learned}).

\paragraph{A held-out judge retains only part of the gain, and how much depends on the
optimizer.} We re-scored every conversation in this comparison --- $16{,}896$ cells, all
eight instruments $\times$ 22 model states $\times$ 96 conversations --- with \heldoutjudge{}, a
grader from a different model family that never played the patient and never participated in
training. We report \emph{gain retention}, $\Delta_\text{held-out}/\Delta_\text{trained-against}$,
as a per-arm train/test generalization ratio. At the matched ten-iteration endpoint, PTO
retains $0.80$ of its Q1 gain and GRPO $0.28$, with non-overlapping bootstrap intervals.
Crucially retention is an \emph{onset} curve, not an endpoint accident: the arms are
indistinguishable for three iterations, then separate --- the shape a policy drifting onto
grader-specific features would produce, legible four iterations before the reward curve turns
over (\S\ref{sec:heldout}).

\paragraph{Contributions.} (i) A controlled ten-iteration comparison of two exploration
strategies under an identical LLM-questionnaire reward in a clinical dialogue task, with all
eight instruments reported flat. (ii) \emph{Gain retention} against a decoupled held-out
judge as a cheap, human-free diagnostic for reward hacking, together with evidence that it is
more informative than the outcome curve and more reliable than rate-based MI-inconsistency
measures. (iii) A behavioural characterisation of what the hack \emph{is} in this domain ---
affirmation without inquiry --- including a grader-versus-regex divergence that localises
part of the failure in the measuring instrument rather than only in the policy. We also show
(\S\ref{sec:discussion}, Appendix~\ref{app:probe}) that the gap between the two optimizers is
attributable to the state distribution their training data is drawn from --- on-policy slices
versus a best-of-$M$ reranked trunk --- rather than to the loss that weights the candidates.
